\begin{ex}
Tính: $\frac{2}{7} + \frac{3}{7} = \dots$
\shortans{5/7}
\loigiai{
$\frac{5}{7}$.
}
\end{ex}

\begin{ex}
Tính: $\frac{7}{9} - \frac{4}{9} = \dots$ (tối giản).
\shortans{1/3}
\loigiai{
$\frac{3}{9} = \frac{1}{3}$.
}
\end{ex}

\begin{ex}
Tính: $\frac{1}{2} + \frac{1}{4} = \dots$
\shortans{3/4}
\loigiai{
$\frac{3}{4}$.
}
\end{ex}

\begin{ex}
Tính: $\frac{3}{4} - \frac{1}{3} = \dots$
\shortans{5/12}
\loigiai{
$\frac{9}{12} - \frac{4}{12} = \frac{5}{12}$.
}
\end{ex}

\begin{ex}
Tính: $\frac{2}{3} \times \frac{4}{5} = \dots$
\shortans{8/15}
\loigiai{
$\frac{8}{15}$.
}
\end{ex}

\begin{ex}
Tính: $\frac{3}{5} : \frac{2}{3} = \dots$
\shortans{9/10}
\loigiai{
$\frac{3}{5} \times \frac{3}{2} = \frac{9}{10}$.
}
\end{ex}

\begin{ex}
Tính: $\frac{4}{7} \times 14 = \dots$
\shortans{8}
\loigiai{
$14 : 7 \times 4 = 8$.
}
\end{ex}

\begin{ex}
Tính: $12 : \frac{3}{4} = \dots$
\shortans{16}
\loigiai{
$12 \times \frac{4}{3} = 16$.
}
\end{ex}

\begin{ex}
Tính: $1 - \frac{3}{8} = \dots$
\shortans{5/8}
\loigiai{
$\frac{5}{8}$.
}
\end{ex}

\begin{ex}
Tính: $1 + \frac{2}{5} = \dots$
\shortans{7/5}
\loigiai{
$\frac{7}{5}$.
}
\end{ex}

\begin{ex}
Tìm $\frac{3}{4}$ của 36.
\shortans{27}
\loigiai{
$36 \times \frac{3}{4} = 27$.
}
\end{ex}

\begin{ex}
Tìm $\frac{2}{5}$ của 50 kg theo đơn vị ki-lô-gam.
\shortans{20}
\loigiai{
$50 \times \frac{2}{5} = 20$ kg.
}
\end{ex}

\begin{ex}
Một tấm bìa hình chữ nhật có chiều dài $\frac{4}{5}$ m, chiều rộng $\frac{1}{2}$ m. Diện tích tấm bìa là bao nhiêu mét vuông? (tối giản).
\shortans{2/5}
\loigiai{
$\frac{4}{5} \times \frac{1}{2} = \frac{2}{5}\text{ m}^2$.
}
\end{ex}

\begin{ex}
Chu vi tấm bìa trong bài toán trên là bao nhiêu mét?
\shortans{13/5}
\loigiai{
$(\frac{4}{5} + \frac{1}{2}) \times 2 = \frac{13}{10} \times 2 = \frac{13}{5}$ m.
}
\end{ex}

\begin{ex}
Tính nhanh: $\frac{2}{5} \times \frac{3}{7} + \frac{2}{5} \times \frac{4}{7}$.
\shortans{2/5}
\loigiai{
$\frac{2}{5} \times 1 = \frac{2}{5}$.
}
\end{ex}

\begin{ex}
Tính giá trị biểu thức: $(\frac{1}{2} + \frac{1}{4}) \times 4$.
\shortans{3}
\loigiai{
$\frac{3}{4} \times 4 = 3$.
}
\end{ex}

\begin{ex}
Tính giá trị biểu thức: $1 - \frac{1}{4} : \frac{1}{2}$.
\shortans{1/2}
\loigiai{
$1 - \frac{1}{2} = \frac{1}{2}$.
}
\end{ex}

\begin{ex}
Tìm $x$: $x + \frac{1}{3} = \frac{5}{6}$.
\shortans{1/2}
\loigiai{
$x = \frac{5}{6} - \frac{2}{6} = \frac{3}{6} = \frac{1}{2}$.
}
\end{ex}

\begin{ex}
Tìm $x$: $x \times \frac{3}{4} = 9$.
\shortans{12}
\loigiai{
$x = 9 : \frac{3}{4} = 12$.
}
\end{ex}

\begin{ex}
Tìm $x$: $x : \frac{2}{5} = 15$.
\shortans{6}
\loigiai{
$x = 15 \times \frac{2}{5} = 6$.
}
\end{ex}

\begin{ex}
Một lớp có 36 học sinh, trong đó $\frac{4}{9}$ là học sinh xuất sắc. Có bao nhiêu học sinh xuất sắc?
\shortans{16}
\loigiai{
$36 \times \frac{4}{9} = 16$ bạn.
}
\end{ex}

\begin{ex}
Tính: $\frac{3}{7} : \frac{3}{7} = \dots$
\shortans{1}
\loigiai{
1.
}
\end{ex}

\begin{ex}
Tính: $\frac{5}{9} \times \frac{9}{5} = \dots$
\shortans{1}
\loigiai{
1.
}
\end{ex}

\begin{ex}
Tính: $\frac{7}{8} - \frac{7}{8} = \dots$
\shortans{0}
\loigiai{
0.
}
\end{ex}

\begin{ex}
Tính: $\frac{1}{2} \times \frac{2}{3} \times \frac{3}{4} \times \frac{4}{5} = \dots$
\shortans{1/5}
\loigiai{
Triệt tiêu được $\frac{1}{5}$.
}
\end{ex}
